\kbvexercise{ಅಭ್ಯಾಸ ಪತ್ರಿಕೆ ೧೭}

\vspace{2.5pt plus 3pt minus 0.5pt}\noindent\hspace*{1.2em}\kbvnum{(1)}\hspace{0.55em}ದ್ವಿರುಕ್ತಿ ಎಂದರೇನು ? ಉದಾಹರಣೆಗಳ ಸಹಿತ ತಿಳಿಸಿ,\par

\vspace{2.5pt plus 3pt minus 0.5pt}\noindent\hspace*{1.2em}\kbvnum{(2)}\hspace{0.55em}ಕೆಳಗೆ ಕೊಟ್ಟಿರುವ ದ್ವಿರುಕ್ತಿಗಳನ್ನು ನಿಮ್ಮ ಸ್ವಂತ ಮಾತುಗಳು\par

\indent (ವಾಕ್ಯಗಳಲ್ಲಿ) ಬಳಸಿರಿ :\par

\indent ಬೇಗಬೇಗ ಈಗೀಗ ಹೆಚ್ಚುಹೆಚ್ಚು ನಡೆನಡೆ\par

\indent ನಿಲುನಿಲು ಮನೆಮನೆ  ಕಮ್ಮಿಕಮ್ಮಿ ಭಲೇಭಲೇ\par

\vspace{2.5pt plus 3pt minus 0.5pt}\noindent\hspace*{1.2em}\kbvnum{(3)}\hspace{0.55em}ಭಾವಸೂಚಕಾವ್ಯಯಗಳನ್ನು ದ್ವಿರುಕ್ತಿಗಳಾಗಿ ಬಳಸಬಹುದೇ?\par

\vspace{2.5pt plus 3pt minus 0.5pt}\noindent\hspace*{1.2em}\kbvnum{(4)}\hspace{0.55em}“ಠೂ ! ಠೂ ! ನೀನೊಬ್ಬ ಅಯೋಗ್ಯ”\par

\indent —ಇಲ್ಲಿ ಠೂ ! ಠೂ ! ಎಂಬುದು ದ್ವಿರುಕ್ತಿಯೇ ? ಅಥವಾ ಅವ್ಯಯವೇ ?\par

\vspace{2.5pt plus 3pt minus 0.5pt}\noindent\hspace*{1.2em}\kbvnum{(5)}\hspace{0.55em}“ಸತ್ಯವೇ ತಾಯಿ ತಂದೆ; ಸತ್ಯವೇ ಬಂಧು ಬಳಗ”—ಇಲ್ಲಿ ಸತ್ಯವೇ ಎಂಬ ಶಬ್ದವು ಪುನರಾವರ್ತನೆಯಾಗಿದೆ. ಇದು ದ್ವಿರುಕ್ತಿಯೇ?\par

\kbvchapter{ಪಾಠ ೧೮}{ತತ್ಸಮ-ತದ್ಭವ}

\indent ಕನ್ನಡ ಭಾಷೆಯಲ್ಲಿ ಪದ ಸಂಪತ್ತು ಬರುಬರುತ್ತಾ ಹೆಚ್ಚುತ್ತಿದೆ. ಅನ್ಯ ಭಾಷಾಪದಗಳ ಕನ್ನಡೀಕರಣ ಈ ಬೆಳವಣಿಗೆಗೆ ಬಹು ಮುಖ್ಯ ಕಾರಣ, ಆಂಗ್ಲ ಭಾಷೆಯ ಪದಗಳು, ಫಾರಸಿ ಭಾಷೆಯ ಪದಗಳು, ಸಂಸ್ಕೃತ ಪದಗಳು ಮತ್ತು ಇನ್ನಿತರ ಭಾರತೀಯ ಭಾಷೆಗಳಿಂದ ಬಂದ ಪದಗಳು ಕನ್ನಡದೊಂದಿಗೆ ಸೇರಿಕೊಂಡು ಕನ್ನಡ ಪದಗಳೇ ಆಗಿಬಿಟ್ಟಿವೆ. ಅವುಗಳಲ್ಲಿ ಸಂಸ್ಕೃತ ಪದಗಳ ಮಾತ್ರ ಹೆಚ್ಚು ಸಂಖ್ಯೆಯಲ್ಲಿವೆ,\par

\indent ಅಪ್ಪಟ ಸಂಸ್ಕೃತ ಪದಗಳನ್ನೂ ಒಂದಿಷ್ಟು ಮಾರ್ಪಾಡು ಹೊಂದಿರುವ ಸಂಸ್ಕೃತ ಪದಗಳನ್ನೂ ಕನ್ನಡದಲ್ಲಿ ಕಾಣಬಹುದು. ಯಾವುದೊಂದು ಬದ ಲಾವಣೆಯೂ ಆಗದೇ ಕನ್ನಡದಲ್ಲಿ ಬಳಕೆಗೆ ಬಂದಿರುವ ಸಂಸ್ಕೃತ ಪದಗಳಿಗೆ “ತತ್ಸಮ”ಗಳೆಂದು ಹೆಸರು; ಸಾಕಷ್ಟು ಮಟ್ಟಿಗೆ ಬದಲಾವಣೆ ಹೊಂದಿ ಕನ್ನಡ ಪದಗಳ ಗುಂಪಿಗೆ ಸೇರಿರುವ ಸಂಸ್ಕೃತ ಪದಗಳಿಗೆ “ತದ್ಭವ”ಗಳೆಂದು ಹೆಸರು.\par

\indent ಉದಾಹರಣೆಗೆ ಸ್ತ್ರೀ-ಅಪ್ಪಟ ಸಂಸ್ಕೃತ ಪದ. ಇದನ್ನು ಕನ್ನಡದಲ್ಲಿ ನೇರವಾಗಿ ಬಳಸುತ್ತೇವೆ; ಇದೊಂದು ತತ್ಸಮ ಶಬ್ದ. ಕರ್ತರಿ-ಅಪ್ಪಟ ಸಂಸ್ಕೃತ\kbvcontinues
